\documentclass[conference]{IEEEtran}
\IEEEoverridecommandlockouts

\usepackage{cite}
\usepackage{amsmath,amssymb,amsfonts}
\usepackage{mathtools}
\usepackage{graphicx}
\usepackage[utf8]{inputenc}
\usepackage[brazilian]{babel}
\usepackage[hidelinks]{hyperref}
\usepackage{textcomp}
\usepackage{accsupp}
\usepackage[T1]{fontenc}
\usepackage{amsmath}            
\usepackage{float}              
\usepackage{algorithm}           
\usepackage{algpseudocode}       
\usepackage{listings}
\usepackage{xcolor}
\def\BibTeX{{\rm B\kern-.05em{\sc i\kern-.025em b}\kern-.08em
    T\kern-.1667em\lower.7ex\hbox{E}\kern-.125emX}}
% --- Anexo de uso de IAG ---
\newcounter{usoiag}
\newcommand{\usoIAG}[4]{%
  \refstepcounter{usoiag}%
  \par\medskip\noindent\hrule\smallskip
  \noindent\textbf{Uso \theusoiag}\\
  \textbf{Etapa / local de uso:} #1\\
  \textbf{Modelo utilizado:} #2\\
  \textit{\textbf{Prompt:}} \textit{#3}\\
  \textbf{Observações:} #4\par
}

\floatname{algorithm}{Algoritmo}
\algrenewcommand\algorithmicrequire{\textbf{Entrada:}}
\algrenewcommand\algorithmicensure{\textbf{Saída:}}
\algrenewcommand\algorithmicfor{\textbf{para}}
\algrenewcommand\algorithmicdo{\textbf{faça}}
\algrenewcommand\algorithmicif{\textbf{se}}
\algrenewcommand\algorithmicthen{\textbf{então}}
\algrenewcommand\algorithmicelse{\textbf{senão}}
\algrenewcommand\algorithmicreturn{\textbf{retorne}}


\definecolor{codebg}{RGB}{245, 245, 245}
\definecolor{codegreen}{RGB}{0, 128, 0}
\definecolor{codegray}{RGB}{128, 128, 128}
\definecolor{codepurple}{RGB}{128, 0, 128}
\definecolor{codeblue}{RGB}{0, 0, 180}


\newcommand*{\nonselectablenumber}[1]{{\BeginAccSupp{ActualText={}}#1\EndAccSupp{}}}


\lstdefinestyle{codigo}{
    backgroundcolor=\color{codebg},
    commentstyle=\color{codegreen},
    keywordstyle=\color{codeblue},
    numberstyle=\tiny\color{codegray}\nonselectablenumber,
    stringstyle=\color{codepurple},
    basicstyle=\ttfamily\scriptsize,
    breakatwhitespace=false,
    breaklines=true,
    captionpos=b,
    nolol=false,
    keepspaces=true,
    numbers=left,
    numbersep=5pt,
    showspaces=false,
    showstringspaces=false,
    showtabs=false,
    tabsize=3,
    frame=single,
    framerule=0.5pt,
    rulecolor=\color{codegray},
    extendedchars=true,
    literate=
    {á}{{\'a}}1 {ã}{{\~a}}1 {â}{{\^a}}1 {à}{{\`a}}1
    {é}{{\'e}}1 {ê}{{\^e}}1 {è}{{\`e}}1
    {í}{{\'i}}1 {î}{{\^i}}1 {ì}{{\`i}}1
    {ó}{{\'o}}1 {ô}{{\^o}}1 {õ}{{\~o}}1 {ò}{{\`o}}1
    {ú}{{\'u}}1 {ü}{{\"u}}1 {ù}{{\`u}}1
    {ç}{{\c{c}}}1 {Ç}{{\c{C}}}1
}

\lstset{style=codigo}

\lstset{style=codigo}

    \begin{document}

\title{Algoritmo Dígito por Dígito para Cálculo de Raiz Quadrada no ATmega328}
\author{
\IEEEauthorblockN{Gabriel Luduvice Nunes Wolfovitch, João Pedro Nascimento Monteiro e \\
João Vítor de Araújo Costa de Oliveira}
\IEEEauthorblockA{\textit{Departamento de Engenharia Elétrica e Computação (DEEC)}\\
\textit{Universidade Federal da Bahia (UFBA)}\\
Salvador, Brasil\\
\{gabrielwolfovitch, joaovaco, email3\}@ufba.br
}
}

\maketitle

\begin{abstract}
O problema proposto envolve o desenvolvimento do cálculo de uma raiz em um microcontrolador ATmega328. Como esse microcontrolador não possui uma unidade de ponto flutuante e há a impossibilidade de utilizar bibliotecas matemáticas vistas em C, este relatório apresenta um algoritmo que calcula a raiz quadrada binária pelo método dígito por dígito, feito inteiramente na linguagem Assembly. O método utiliza representação em ponto fixo, processando uma entrada de 16 bits (Q16.0) para gerar uma saída (Q8.6). A validação foi feita por meio de um teste prévio em C e testes de rastreamento, garantindo que o método dígito por dígito gera uma execução eficiente, ideal para o cálculo de variáveis.
\end{abstract}

\begin{IEEEkeywords}
Raiz Quadrada, Assembly, Linguagem C, ATmega328.
\end{IEEEkeywords}

\section{Introdução}
O processamento em um sistema embarcado de instrumentação baseado em microcontroladores necessita realizar operações matemáticas básicas sobre as grandezas medidas, a exemplo da radiciação. Isto é, durante certas medições, o sistema recebe de sensores grandezas proporcionais ao quadrado de uma variável física. Para recuperar esta variável, é preciso calcular a raiz quadrada do valor medido.

No entanto, dispositivos como o ATmega328 não possuem funções nativas para o cálculo de raízes quadradas, necessitando implementações diretas. Frente às restrições relacionadas ao projeto, não é possível fazer uso de bibliotecas matemáticas, nem utilizar ponto flutuante.

Diante disso, a implementação via dígito por dígito faz uso de um algoritmo simples, baseado apenas em deslocamentos, somas e subtrações entre registradores, para realizar a operação desejada. Por se basear em agrupamentos e subdivisões, os erros associados são tão baixos quanto menor for a quantização em ponto fixo. O repositório com o código da implementação final pode ser encontrado em \cite{monteiro_sismicro}.

\section{Justificativa do Dígito por Dígito}
A escolha da extração de raiz quadrada binária dígito a dígito foi baseada na sua simplicidade. Considerando operações em ponto fixo, este método é preferível por sua facilidade no acompanhamento e checagem, além de não necessitar de operações complexas.

Toda a lógica matemática é resolvida apenas através de:
\begin{itemize}
    \item \textbf{Deslocamentos lógicos:} Operações já nativas das instruções do ATmega328, não necessitando grandes linhas de código para sua construção.
    \item \textbf{Subtrações:} Implementadas também nativamente entre os registradores com as instruções existentes.
    \item \textbf{Testes de Sinal:} Avaliação direta do bit de sinal do SREG para tomadas de decisão.
\end{itemize}

Dessa forma, o custo computacional passa a ser linear em relação ao número de bits do radicando. Além disso, não há necessidade de aramzenar valores em uma \textit{Lookup Table} (LUT), nem de realizar multiplicações ou divisões.

\section{Funcionamento do Algoritmo}
Em linhas gerais, o algoritmo consiste em subdividir o número binário em subgrupos de 2 dígitos. A cada iteração, cada um desses grupos irá gerar um único dígito como resposta. Como a base binária permite apenas 0 ou 1, o dígito seguinte é decidido através da diferença entre o cálculo restante e o valor previsto.

Detalhes acerca da demonstração matemática, bem como outros exemplos, podem ser encontrados em \cite{wiki_sqrt}. O algoritmo abaixo generaliza o exemplo mostrado em \cite{projectf_sqrt}
\begin{algorithm}[htbp]
\caption{Raiz quadrada inteira binária em ponto fixo}
\label{alg:sqrt-registradores}
\begin{algorithmic}[1]
\Require Radicando $X$ de $2n$ bits (sem sinal)
\Ensure Raiz $Q$ de $n$ bits

\State $A \gets 0$ \Comment{acumulador}
\State $Q \gets 0$ \Comment{raiz parcial}
\For{$i \gets 1$ \textbf{até} $n$}
    \State $A \gets (A \ll 2) \mid X[2n\!-\!1\!:\!2n\!-\!2]$ \Comment{2 bits de $X$ entram em $A$}
    \State $X \gets X \ll 2$
    \State $T \gets A - \{Q,\,01_2\}$ \Comment{$\{Q,01_2\} = (Q \ll 2) \mid 01_2$}
    \State $Q \gets Q \ll 1$
    \If{$T \geq 0$}
        \State $A \gets T$
        \State $Q[0] \gets 1$
    \EndIf
\EndFor
\State \Return $Q$
\end{algorithmic}
\end{algorithm}

\newpage
\section{Modelagem em Alto Nível: Linguagem C}
A primeira etapa do projeto consistiu em validar o funcionamento do algoritmo de forma simplificada e rápida. Para isso, realizou-se uma implementação na linguagem C. O código abaixo demonstra a lógica principal.

\begin{lstlisting}[language=C, title={Implementação em C do algoritmo dígito por dígito.}, label={lst:codigoC}]
    #include <stdint.h>
    #include <stdio.h>

    typedef uint8_t REG_t;

    typedef struct ufix16 {
    REG_t MSB;
    REG_t LSB;
    } ufix16_t;

    float evaluate_ufix16_q88(const ufix16_t *num);
    float evaluate_ufix16_q160(const ufix16_t *num);
    ufix16_t take_sqrt(const ufix16_t *input);

    int main(void) {
    // assume input Q16.0 e output Q8.8
    ufix16_t num = {0x0, 0xff};
    printf("input is %f\n", evaluate_ufix16_q160(&num));
    ufix16_t res = take_sqrt(&num);
    printf("result is %f\n", evaluate_ufix16_q88(&res));

    return 0;
    }

    float evaluate_ufix16_q88(const ufix16_t *num) {
    uint16_t raw_value = ((uint16_t)num->MSB << 8) | num->LSB;
    return (float)raw_value / 256.0f;
    }

    float evaluate_ufix16_q160(const ufix16_t *num) {
    uint16_t raw_value = ((uint16_t)num->MSB << 8) | num->LSB;
    return (float)raw_value;
    }

    ufix16_t take_sqrt(const ufix16_t *input) {
    // Alinha a entrada de 16 bits (Q16.0) em um espaço de 32 bits.
    // Isso efetivamente multiplica o número por 2^16, convertendo para Q16.16.
    uint32_t val = ((uint32_t)input->MSB << 24) | ((uint32_t)input->LSB << 16);

    uint32_t res = 0; // Acumulador da raiz quadrada (resultado)

    // O 'bit' de teste começa na maior potência de 4 que cabe em 32 bits.
    uint32_t bit = 1UL << 30;
    // Alinha o 'bit' inicial ao tamanho real do valor de entrada
    while (bit > val) {
        bit >>= 2;
    }
    // Laço de extração bit a bit
    while (bit != 0) {
        if (val >= res + bit) {
        val = val - (res + bit);
        res = (res >> 1) + bit;
        } else {
        res >>= 1;
        }
        bit >>= 2;
    }
    // 'res' agora contém a raiz exata no formato Q8.8
    ufix16_t result;
    result.MSB = (res >> 8) & 0xFF;
    result.LSB = res & 0xFF;
    return result;
}
\end{lstlisting}

A execução deste código em C auxiliou na compreensão mais abstrata do funcionamento geral do projeto. Assim, este foi adotado como base para a implementação final em Assembly.

\section{Descrição do Sistema Final}
\subsection{Quantizações e Registradores}
Para permitir uma maior amplitude de valores numéricos, utilizou-se 2 registradores para cada um dos termos descritos no algoritmo \ref{alg:sqrt-registradores}, além de 2 adicionais para a construção do valor $\{Q,\,01_2\}$. Suas identificações e faixas de endereço são:
\begin{enumerate}
    \item \textbf{Entrada X}  (R19:18);
    
    \item \textbf{Resposta Q} (R21:R20);
    
    \item \textbf{Valor sendo trabalhado A} (R23:R22);
 
    \item \textbf{Registrador do sinal T} (R25:R24);

    \item \textbf{Resposta temporária $\{Q,\,01_2\}$} (R27:R26).

\end{enumerate}

O sistema final pressupõe uma entrada já quantizada em Q16.0. A fim de cobrir uma ampla faixa dos números inteiros, decidimos interpretar a entrada como um número de 16 bits, dos quais todos são parte inteira. Assim, o código é capaz de calcular a raiz de um número $n$ na forma: $\forall n \in \mathbb{Z}, n \in [0, 65535]$. 

Além disso, a saída foi quantizada em Q8.6. O tamanho da parte inteira da saída, conforme mostrado no algoritmo \ref{alg:sqrt-registradores}, depende apenas do tamanho da entrada. A parte fracionária, para os casos de raízes não exatas, está relacionada com a quantidade de iterações realizadas após preenchimento da parte inteira.

Como cada iteração gera um dígito, o preenchimento da parte fracionária foi feito até o limite do \textit{overflow} dos registradores intermediários do cálculo.

\subsection{Rotinas Principais}
O código Assembly final tem um \textit{loop} com duas rotinas principais, a saber:
\begin{enumerate}
    \item SHIFT\_DOIS;
    \item CALCULA\_SINAL.
\end{enumerate}

Na primeira, 2 bits da entrada são carregado para os registradores que irão guardar o valor sendo atualmente trabalhado. Este processo é feito através de instruções com LSL e ROL, aproveitando o carry que possa ter surgido do primeiro deslocamento. Ambos são executados duas vezes em \textit{loop}. 

Na segunda etapa, ocorre a validação do sinal e a construção da resposta. O resultado atual é salvo em outro par de registradores através do MOVW, deslocado duas vezes e somado 1. Ou seja, este é o registrador $\{Q,\,01_2\}$. 

Em seguida, por meio de outro MOVW e subtrações sem e com carry (SUB e SBC), a resposta será comparada por meio dos registradores A e T. O controle dessa condicional é realizado pelo bit N do SREG, indicando se houve ou não um valor negativo durante a operação. Este processo é repetido até atingir uma saída na formatação escolhida. 

O diagrama completo do sistema final pode ser encontrado em \cite{monteiro_sismicro}.

\newpage
\section{Análise dos Resultados}
Para validar o funcionamento do código em Assembly, realizou-se um teste com alguns casos particulares:
\begin{itemize}
    \item Extremos: Valores no limite inferior e superior do intervalo possível de representação para números em ponto fixo Q16.0
    \item Raízes Exatas: Números que possuem raízes inteiras (``quadrados perfeitos'')
    \item Raízes Inexatas: Números que não possuem raízes inteiras.
\end{itemize}

A Tabela \ref{tabela_teste} traz o resultado de alguns desses casos, mostrando qual foi a saída gerada a partir de cada entrada.

\begin{table}[htbp]
\centering
\caption{\textsc{Casos de Teste da Raiz Quadrada (Entrada Q16, Saída Q8.6)}}
\label{tabela_teste}
\begin{tabular}{lcc}
\hline
\textbf{Entrada em Q16.0} & \textbf{Saída em Q8.6} \\
\hline
0000000001111001 \ (121)   & 00001011.000000 \ (11{,}000000) \\
0000000000000000 \ (0)     & 00000000.000000 \ (0{,}000000) \\
1111111111111111 \ (65535) & 11111111.111111 \ (255{,}984375) \\
0000000000000010 \ (2)     & 00000001.011010 \ (1{,}406250) \\
\hline
\end{tabular}
\end{table}

Para os casos em questão, podemos perceber que o algoritmo cumpriu com o esperado. Para os casos em que a entrada possuía raiz inteira, o cálculo foi exato, sem qualquer erro envolvido. Para os casos onde a raiz não era exata, pudemos comprovar que o erro esteve sempre limitado ao menor valor passível de representação ($2^{-6} \approx 0,015625$).

Isto pode ser mostrado olhando para a entrada máxima possível no intervalo de representação. O valor é $(0\text{x}FFFF)_{16}  = (65535)_{10}$. A raiz então seria $\sqrt{65535} \approx 255,998047$. 

Seja $\delta$ o erro absoluto associado. Ficamos com:
\begin{align*}
    & \delta = 255,998047 - 255,984375 = 0.013672 \\
    & \delta < 2^{-6} = 0,015625
\end{align*}

Conforme era esperado, já que não há ocorrência de \textit{overflow} no processo.

\section{Conclusão}
Este trabalho demonstrou sucesso na implementação de um algoritmo de baixo nível para um microcontrolador. A escolha do método do dígito por dígito se mostrou uma boa abordagem para a arquitetura do ATmega328, extraindo o melhor possível do método de ponto fixo. A adoção do ponto fixo permitiu não apenas processar completamente os inteiros de 16 bits (Q16.0), mas também garantir precisão fracionária na resposta (Q8.6), chegando no limite de \textit{overflow} dos registradores. Como mostrado nos testes de rastreamento, o controle de fluxo por meio da \textit{flag} de sinal do registrador de status e o mapeamento direto das variáveis para os registradores (R18 a R27) garantiram um \textit{datapath} enxuto. Pode-se concluir que o código desenvolvido em Assembly atende aos requisitos solicitados, oferecendo um tempo de execução curto e um baixo uso de memória.

\appendix[Uso de Inteligência Artificial Generativa]

\subsection*{Declaração}
A equipe declara que a IAG foi utilizada nas etapas descritas abaixo. O detalhamento de cada uso, incluindo os \textit{prompts} empregados, é apresentado a seguir.

\subsection*{Detalhamento}

\usoIAG
  {Generalização do código em Assembly}
  {\textit{Claude Sonnet} 5.5 \textit{(Anthropic)}}
  {como esse código poderia ser escalado para lidar com números em uma faixa maior de representação? Por exemplo, em Q8.8?. Dê um direcionamento e ideias para uma implementação, sem enviar o código pronto.}
  {O \textit{prompt} foi feito após já existir uma implementação funcional do algoritmo, porém para uma escala menor de representação, com apenas um registrador.}

\usoIAG
  {Formatação das referências bibliográficas}
  {\textit{Claude Sonnet} 5.5 \textit{(Anthropic)}}
  {formate os links de referências abaixo em um ambiente thebibliography para latex
   em conformidade com o template IEEE conferences em português}
   {O \textit{prompt} acompanhava a lista de referências mostrada ao final do relatório.}

\usoIAG
  {Generalização do código em Assembly}
  {\textit{Claude Sonnet} 5.5 \textit{(Anthropic)}}
  {como esse código poderia ser escalado para lidar com números em uma faixa maior de representação? Por exemplo, em Q8.8?. Dê um direcionamento e ideias para uma implementação, sem enviar o código pronto.}
  {O \textit{prompt} foi feito após já existir uma implementação funcional do algoritmo, porém para uma escala menor de representação, com apenas um registrador para a resposta. Acompanhava também o código existente até então.}

\usoIAG{Formatação e generalização do algoritmo para relatório}
        {\textit{Claude Sonnet} 5.5 \textit{(Anthropic)}}
        {escreva isso genericamente em forma de um bloco latex de algoritmo para relatório}
        {Adaptado algumas vezes para adequar nomenclaturas. O \textit{prompt} acompanhava a imagem de um exemplo presente em \cite{projectf_sqrt}.}

\usoIAG
  {Formatação da seção anexa de IAG}
  {\textit{Claude Sonnet} 5.5 \textit{(Anthropic)}}
  {faça um template para que eu possa colocar essa seção como um anexo do arquivo já existente no latex no template IEEE Conference. Alguns campos para o prompt, indicação de onde foi usado e o modelo usado}
   {O \textit{prompt} incluía uma imagem com as orientações para a seção de IAG. Adaptado em relação à formatação.}


\begin{thebibliography}{1}
\bibitem{monteiro_sismicro}
``sismicro: Unidade 1,'' repositório GitHub. [Online]. Disponível em: \url{https://github.com/JoaoNasMonteiro/sismicro/tree/main/u1}. Acesso em: 18 out. 2026.

\bibitem{wiki_sqrt}
``Square root algorithms: Digit By Digit Calculation,'' \emph{Wikipedia, The Free Encyclopedia}. [Online]. Disponível em: \url{https://en.wikipedia.org/wiki/Square_root_algorithms#Digit-by-digit_calculation}. Acesso em: 18 out. 2026.

\bibitem{projectf_sqrt}
W. Green, ``Square root in Verilog,'' \emph{Project F}, 22 dez. 2020. [Online]. Disponível em: \url{https://projectf.io/posts/square-root-in-verilog/}. Acesso em: 18 out. 2026.

\bibitem{morris_goldschmidt}
M. Morris, ``Computing fixed-point square roots and their reciprocals using Goldschmidt algorithm,'' \emph{FPGARelated.com}, 14 jun. 2020. [Online]. Disponível em: \url{https://www.fpgarelated.com/showarticle/1347.php}. Acesso em: 18 out. 2026.

\bibitem{ncot_fixedpoint}
NCOT Technology, ``Fixed point maths explained - Retro programming,'' YouTube, 14 nov. 2021. [Vídeo]. Disponível em: \url{https://www.youtube.com/watch?v=Is67DfCdvcE}. Acesso em: 18 out. 2026.
\end{thebibliography}

\end{document}